\chapter{The Pro-etale Site as a Cohomological Workspace}
\label{chap:proetale-site}

The pro-etale site is not a replacement for the etale topology in the
finiteness theorems.  Its advantage is different and very concrete: inverse
limits of finite etale covers become actual objects, complete topological
coefficient rings become honest sheaves, and derived inverse limits can be
controlled internally.  This chapter isolates precisely those gains.

Throughout, schemes are locally noetherian unless explicitly stated.  The
definitions themselves make sense for arbitrary schemes.

\section{Weakly etale morphisms}

\begin{definition}[Weakly etale morphism]
\label{def:weakly-etale}
A morphism \(f:Y\to X\) is \emph{weakly etale} if \(f\) is flat and its
diagonal
\[
  \Delta_f:Y\longrightarrow Y\times_XY
\]
is flat.  It is \emph{pro-etale} in the affine sense if, locally on source
and target, it is a cofiltered inverse limit of etale affine morphisms with
affine transition maps.
\end{definition}

Every etale morphism is weakly etale.  A filtered colimit of etale algebras
is weakly etale, so inverse limits of affine etale schemes provide the basic
new examples.  Conversely, weakly etale morphisms are locally ind-etale in
the topology relevant below; this is
\cite[Theorem~2.3.4]{BhattScholzeProEtale}.

\begin{example}[A profinite set becomes geometric]
Let \(S=\varprojlim_i S_i\) be a profinite set, written as a limit of finite
sets, and let \(X\) be a scheme.  Then
\[
  X\times S:=\varprojlim_i \coprod_{s\in S_i}X
\]
is pro-etale over \(X\).  Thus a continuous action of a profinite group may
be studied on an actual cover rather than only through all of its finite
quotients.
\end{example}

\begin{definition}[Small pro-etale site]
\label{def:small-proetale-site}
The site \(X_{\proet}\) has weakly etale \(X\)-schemes as objects.  Coverings
are generated by jointly surjective fpqc families.  Equivalently, one may
use the basis of affine pro-etale objects and finite jointly surjective
families after refining.  We write
\[
  \nu_X:X_{\proet}\longrightarrow X_{\et}
\]
for the morphism of topoi induced by viewing every etale object as a
pro-etale object.
\end{definition}

The equivalence of the weakly-etale and ind-etale descriptions of the
topology is a theorem, not a definitional equality.  This distinction matters
in a formal development: morphisms, covering families, and the chosen basis
should be separate structures.

\begin{proposition}[Full faithfulness for bounded-below etale complexes]
\label{prop:proetale-etale-comparison}
For \(K\in\D^+(X_{\et})\), the adjunction unit
\[
 K\longrightarrow \R\nu_{X,*}\nu_X^*K
\]
is an equivalence.  Hence
\(\nu_X^*:\D^+(X_{\et})\to\D^+(X_{\proet})\) is fully faithful, and
\[
 \R\Gamma(X_{\et},K)
   \simeq \R\Gamma(X_{\proet},\nu_X^*K).
\]
\end{proposition}

\begin{proof}
If an affine pro-etale object has a presentation \(U=\varprojlim_iU_i\) by
affine etale objects, then
\[
 \R\Gamma(U,\nu_X^*K)
   \simeq\varinjlim_i\R\Gamma(U_i,K).
\]
Refining covers by limits of affine etale covers reduces the assertion to
Cech descent on \(X_{\et}\).  This is precisely
\cite[Corollary~5.1.6]{BhattScholzeProEtale}.  Bounded-below is essential;
\cite[Remark~5.1.8]{BhattScholzeProEtale} gives the unbounded correction by
a Postnikov limit and explains why the displayed unit can then fail.
\end{proof}

\begin{proposition}[Pushforward comparison]
\label{prop:proetale-pushforward-comparison}
Let \(f:Y\to X\) be qcqs and let \(K\in\D^+(Y_{\et})\).  The canonical
exchange morphism
\[
 \nu_X^*\R f_{\et,*}K
   \longrightarrow
 \R f_{\proet,*}\nu_Y^*K
\]
is an equivalence.
\end{proposition}

\begin{proof}
For affine \(X\) and \(Y\), evaluate on a presentation
\(U=\varprojlim_iU_i\) and use the preceding colimit formula on
\(f^{-1}U=\varprojlim_i f^{-1}U_i\).  A finite Mayer--Vietoris argument first
handles quasicompact separated \(Y\), and a second application handles the
qcqs case.  For complexes, finite limits are replaced by finite homotopy
limits.  This is \cite[Lemma~5.4.3]{BhattScholzeProEtale}.
\end{proof}

\begin{warning}
This comparison does not say that every pro-etale sheaf comes from the etale
site.  In particular, the sheaf of continuous \(\Zl\)-valued functions is not
the pullback of the constant discrete sheaf with underlying group \(\Zl\).
That difference is exactly what makes integral \(\ell\)-adic coefficients
work cleanly.
\end{warning}

\section{Weakly contractible affine objects}

\begin{definition}
A qcqs scheme \(U\) is \emph{w-local} if its underlying spectral space
satisfies both:
\begin{enumerate}
\item every open cover \((U_i\hookrightarrow U)_i\) splits, meaning that
      \(\coprod_iU_i\to U\) admits a section;
\item the subspace \(U^c\subset U\) of closed points is closed.
\end{enumerate}
Equivalently, \(U^c\) is closed and the natural map
\[
  U^c\longrightarrow\pi_0(U)
\]
is a bijection, indeed a homeomorphism onto the profinite component space.
An affine w-local scheme is \emph{w-strictly local} if every local ring at a
closed point is strictly henselian.  It is \emph{w-contractible} if every
faithfully flat weakly etale map \(V\to U\) admits a section.
\end{definition}

The equivalence in the definition is
\cite[Definition~2.1.1 and Lemma~2.1.4]{BhattScholzeProEtale}.

The last condition is the site-theoretic one: every covering of the
representable sheaf \(h_U\) splits.  It implies exactness of sections on
epimorphisms of sheaves over \(U\).

\begin{theorem}[Existence of a contractible basis]
\label{thm:proetale-wcontractible-cover}
Every affine scheme admits a faithfully flat ind-etale cover by a
w-contractible affine scheme.  Hence every scheme is locally weakly
contractible on its pro-etale site.
\end{theorem}

\begin{proof}
First pass to a w-local ind-Zariski cover.  Strictly henselize at all closed
points simultaneously, then take the inverse limit over all finite etale
covers equipped with compatible lifts at those points.  The resulting affine
scheme is w-strictly local and has no nontrivial weakly etale covering without
a section.  The construction and the final lifting argument are
\cite[Lemma~2.4.9 and Proposition~4.2.8]{BhattScholzeProEtale}; the global
form is \cite[Theorem~1.5]{BhattScholzeProEtale}.
\end{proof}

\begin{proetalebox}
On the etale site, a geometric point probes a sheaf but is not itself a large
enough neighborhood on which an infinite inverse system becomes constant.
On the pro-etale site, w-contractible affines play the role of projective
test objects.  This turns many derived questions about complete coefficients
into ordinary algebra over the ring of sections on such an affine.
\end{proetalebox}

\section{Repleteness and derived inverse limits}

\begin{definition}
A topos is \emph{replete} if the inverse limit of every countable tower of
epimorphisms is again an epimorphism.
\end{definition}

\begin{theorem}[Homological consequences of the pro-etale basis]
\label{thm:proetale-repleteness}
The pro-etale topos of a scheme is replete.  More precisely:
\begin{enumerate}
\item countable products of abelian sheaves are exact;
\item if \(\cdots\twoheadrightarrow F_2\twoheadrightarrow F_1\) is a tower
      of epimorphisms, then \(\varprojlim F_n\to F_m\) is an epimorphism and
      \(\R\varprojlim F_n\simeq\varprojlim F_n\);
\item for an arbitrary countable tower, \(\R\varprojlim\) has cohomological
      dimension at most one;
\item the unbounded derived category is left complete:
      \[
        K\xrightarrow{\sim}
        \R\varprojlim_n \tau^{\ge -n}K;
      \]
\item complexes satisfy hyperdescent on \(X_{\proet}\).
\end{enumerate}
\end{theorem}

\begin{proof}
Evaluate a tower after pulling it to a w-contractible affine cover.  Since
every epimorphism there splits locally, the elementary set-theoretic lifting
argument for a countable inverse tower applies.  Conservativity of the cover
globalizes the conclusion.  Exact products and the standard two-term
\(1-\mathrm{shift}\) resolution
\[
  \prod_n K_n \xrightarrow{\,1-\mathrm{shift}\,} \prod_n K_n
\]
give the cohomological-dimension bound.  The truncation tower then converges,
which yields left completeness and hyperdescent.  See
\cite[Propositions~3.1.9--3.1.11, 3.3.3, and 3.3.6]{BhattScholzeProEtale}.
\end{proof}

\begin{remark}
Repleteness is not the assertion that inverse limits are exact without
hypotheses.  The only possible obstruction for a countable tower is the
first derived limit, and it vanishes for epimorphic transition maps.  Later
completion arguments will use exactly this statement.
\end{remark}

\section{Honest topological coefficient sheaves}

\begin{definition}[Continuous-function sheaf]
Let \(M\) be a topological abelian group.  On an affine pro-etale object
\(U\), define
\[
  \underline M(U)=\operatorname{Cont}(|U|,M),
\]
and sheafify.  If \(M=\varprojlim_n M_n\) is profinite with finite discrete
quotients, then
\[
  \underline M \simeq \varprojlim_n \underline{M_n}
\]
as sheaves on \(X_{\proet}\).
\end{definition}

\begin{proposition}[The \(\Zl\)-sheaf]
\label{prop:proetale-zl-sheaf}
For \(\ell\) invertible on \(X\), the topological ring \(\Zl\) defines an
honest sheaf of rings
\[
  \underline{\Zl}
     =\varprojlim_n \underline{\mathbf Z/\ell^n\mathbf Z}
\]
on \(X_{\proet}\).  Its restriction to a w-contractible affine \(U\) is the
ring of continuous maps \(|U|\to\Zl\).
\end{proposition}

\begin{proof}
Continuous maps into a profinite limit are precisely compatible systems of
continuous maps into the finite discrete quotients.  Sheaf descent follows
from fpqc descent for locally constant finite functions and repleteness for
the inverse limit.  This is the construction of
\cite[Lemma~4.2.12]{BhattScholzeProEtale}.
\end{proof}

\begin{warning}[Discrete versus continuous constants]
The constant sheaf associated with the discrete group underlying \(\Zl\)
has locally constant sections.  The topological sheaf \(\underline{\Zl}\)
has continuous sections and is the inverse limit of its finite quotients.
They need not agree.  Every later occurrence of \(\Zl\) as a coefficient
ring means the latter sheaf.
\end{warning}

\section{Invariance under universal homeomorphisms}

\begin{theorem}
\label{thm:proetale-universal-homeomorphism}
If \(v:X'\to X\) is a universal homeomorphism, then pullback induces an
equivalence
\[
  v^*:\operatorname{Sh}(X_{\proet})
      \xrightarrow{\sim}
      \operatorname{Sh}(X'_{\proet}).
\]
It is compatible with the comparison to the etale sites and with the
topological coefficient sheaves above.
\end{theorem}

\begin{proof}
A weakly etale morphism descends uniquely across a universal homeomorphism,
and the underlying topological spaces, hence covering families, are
unchanged.  Thus the induced functor of sites is an equivalence.  See
\cite[Lemma~5.4.2]{BhattScholzeProEtale}.
\end{proof}

This theorem is the formal reason geometric Frobenius, although not an
isomorphism of schemes, acts invertibly on the coefficient theory developed
below.

\begin{formalization}
A reusable formalization should expose four layers:
\begin{enumerate}
\item weakly etale morphisms and their stability under composition and base
      change;
\item the pro-etale coverage and the comparison morphism \(\nu_X\);
\item w-contractible affine objects as a basis with a projectivity property;
\item replete topoi and abstract lemmas on \(\R\varprojlim\), independent of
      algebraic geometry.
\end{enumerate}
The last layer belongs in general homological algebra.  The coefficient
chapters should import it, not reprove it for \(\ell\)-adic sheaves.
\end{formalization}

\begin{auditbox}
Nothing in this chapter proves constructibility, finite cohomological
dimension, proper base change, or purity.  The pro-etale site provides the
correct arena for limits and continuous coefficients.  The hard geometric
theorems remain inputs and will be identified explicitly in
\cref{chap:proetale-six}.
\end{auditbox}

